\section{Computational methods}
\label{sec:methods}

\subsection{Data and reproducibility}
We preserved the supplied observations and first reproduced the unchanged
inherited simulations twice. The resulting trajectories were byte-identical
within each dataset. We then exported the teaching notebooks to Python scripts
with cell markers and source hashes, checking equality of their scientific
abstract syntax trees. The delay study uses separate Python modules for data,
simulation, inference, and result generation.

The data audit found 15 CSV files. Two validation files were numeric duplicates
of the sine and step calibration files, respectively. We excluded both duplicates
from validation, leaving 11 files. We fitted the sine experiment, the step
experiment, and their combination separately. Each calibration experiment has
901 temperature observations. The first observation conditions the initial
temperature and is excluded from the fitted residual vector, leaving 900
residuals per experiment. This treatment does not propagate uncertainty in the
initial temperature or in the ambient temperature inferred from it.

Long calculations run from immutable copies of committed source and raw data.
The run manifests retain source, configuration and data hashes, package versions,
random seeds, every optimizer attempt, and failures. A single manifest identifies
each numerical phase. Result tables and figures are generated from saved outputs.

\subsection{Simulation and fitting}
We used exact state propagation of the linear time-invariant dynamics over each
sampling interval, with linear interpolation of the measured heater signal.
Separate checks compared left and right input holds and the inherited backward
Euler discretization. The pure-delay simulator integrates across shifted input
knots. These choices distinguish input reconstruction and numerical discretization
from a change in model equations.

The baseline comparison fitted seven orders, $n\in\{1,2,3,4,6,8,12\}$, two
initial-state assumptions, and both parameter coordinates. Pure-delay fits used
the same two initial-state assumptions and coordinates. The original model used
its four physical parameters. These choices gave 99 calibration cases.
We minimized unweighted temperature sum of squared errors using bounded
least squares in logarithmic parameters, with 24 starts per case and seed 532.
The broad input--output bounds were $0.2\leq K\leq1.5$,
$5\leq\tau\leq600\,\mathrm{s}$, and
$0.01\leq\theta\leq200\,\mathrm{s}$. The physical-coordinate bounds were
$0.032/1.5\leq\Ua\leq0.032/0.2\,\mathrm{W\,K^{-1}}$,
$0.1\leq\Cp\leq100\,\mathrm{J\,K^{-1}}$, and the same delay bounds.
We retained convergence status and active bounds rather than interpreting every
optimizer return as a successful estimate.

\subsection{Local uncertainty and estimability}
For residual vector $\bm e$ and log-parameter vector $\bm\eta$, the sensitivity
matrix $J=\partial\bm e/\partial\bm\eta$ defines the conditional information
matrix $F=J^\top J/\sigma^2$ under independent Gaussian errors of variance
$\sigma^2$. We computed the singular value decomposition of $J$ directly and
reported the eigenvalues and directions of $F$. Numerical rank uses a relative
singular-value cutoff of $10^{-8}$. This threshold is a numerical convention;
the explicit equivalence classes in Appendix~\ref{sec:structural} establish the
structural conclusions independently.

We reported local uncertainty under a fixed hypothetical sensor standard
deviation of $0.25\,{}^\circ\mathrm C$ and under a common variance estimated
from each calibration residual vector. The fixed value is a sensitivity scenario,
not a calibration established by these experiments. Rank-deficient full parameter
covariances are undefined. A diagnostic pseudoinverse is retained separately and
is not used as a confidence interval for a null direction. We also computed
heteroskedasticity and autocorrelation consistent sandwich estimates at lag
bandwidths of 10, 30, and 60 samples, together with residual autocorrelation
functions and Ljung--Box statistics. These calculations assess sensitivity to
serial dependence; they do not correct a biased mean model.

\subsection{Likelihood profiles and resampling}
For each native parameter we fixed its value and reoptimized all nuisance
parameters, following the profile construction of \citet{raue-2009-structural}.
The exploration combined a global grid with local refinement, targeting 61 points
and four nuisance starts at each point. The exact count after merging grids is
recorded per profile. Additional points examined the exchanged
pole mode of the one-stage cold-start model. We retained disconnected accepted
components and endpoints at the exploration bounds. An endpoint at such a bound
does not establish an unbounded mathematical interval.

For $N$ residuals and sum of squared errors $S(\bm\eta)$, the likelihood-ratio
statistics were $[S(\bm\eta)-S(\hat{\bm\eta})]/\sigma^2$ for fixed variance
and $N\log[S(\bm\eta)/S(\hat{\bm\eta})]$ for a common profiled variance.
The nominal 95\% reference threshold was the 0.95 quantile of a chi-squared
distribution with one degree of freedom. These thresholds assume an appropriate
independent Gaussian likelihood and regular identification; they are conditional
diagnostics when those assumptions fail.

The resampling design specified 400 refits for each of four circular residual
block lengths (1, 10, 30, and 60 samples), plus 400 independent Gaussian
simulations at the fixed $0.25\,{}^\circ\mathrm C$ standard deviation, for each
input--output calibration case. Residuals were centered and resampled within
each experiment; first temperatures, ambient estimates, and heater calibration
were fixed. Each draw was transformed exactly into physical coordinates.
Structural null directions were marked undefined rather than assigned an interval
from arbitrary optimizer selections. These resampling distributions remain
conditional on the mean model and on the assumed within-experiment stationarity.

\subsection{Prospective experiment design}
The experiment-design phase compares finite-horizon input candidates using
scaled parameter information and separate model-discrimination criteria.
All designs are simulations. The final design specification, numerical results,
and constraint checks will be inserted after that phase completes; no proposed
waveform has been applied to laboratory hardware.
